% trading-function.tex — Trading function derivation and canonical form verification
% Kontrol: prove_cfmm_replication_tradingFunctionConcavity

\section{Trading Function}\label{sec:trading-function}

We derive the trading function by eliminating the price parameter $p$ from the
reserve curves $x(p)$ and $y(p)$ established in \cref{sec:reserves}.

\subsection{Derivation}

\begin{theorem}[Trading Function]\label{thm:trading-function}
The trading function for the fee concentration derivative CFMM is:
\begin{equation}\label{eq:trading-function-explicit}
  y = x - \ln(x) - 1
\end{equation}
or equivalently in implicit form:
\begin{equation}\label{eq:trading-function-implicit}
  \psi(x, y) = y + \ln(x) + 1 - x = 0
\end{equation}
with domain $x \in (0, 1]$ and $y \in [0, \infty)$.
\end{theorem}

\begin{proof}
From \eqref{eq:risky-reserve}, $x = 1/(1+p)$, so $p = 1/x - 1 = (1-x)/x$.

Substituting into \eqref{eq:numeraire-reserve}:
\begin{align}
  y &= \ln\bigl(1 + \tfrac{1-x}{x}\bigr) - \frac{(1-x)/x}{1 + (1-x)/x} \notag\\
    &= \ln\bigl(\tfrac{1}{x}\bigr) - \frac{(1-x)/x}{1/x} \notag\\
    &= -\ln(x) - (1-x) \notag\\
    &= x - \ln(x) - 1 \label{eq:y-from-x}
\end{align}
Rearranging: $y + \ln(x) + 1 - x = 0$, giving $\psi(x,y) = y + \ln(x) + 1 - x$. \qedhere
\end{proof}

\subsection{Concavity Verification}

\begin{proposition}[Trading Function Concavity]\label{prop:psi-concave}
$\psi(x,y) = y + \ln(x) + 1 - x$ is concave in $(x,y)$ for $x > 0$.
\end{proposition}

\begin{proof}
The Hessian of $\psi$ is:
\begin{equation}\label{eq:hessian}
  H_\psi = \begin{pmatrix}
    \dfrac{\partial^2 \psi}{\partial x^2} & \dfrac{\partial^2 \psi}{\partial x \partial y} \\[8pt]
    \dfrac{\partial^2 \psi}{\partial y \partial x} & \dfrac{\partial^2 \psi}{\partial y^2}
  \end{pmatrix}
  = \begin{pmatrix}
    -1/x^2 & 0 \\
    0 & 0
  \end{pmatrix}
\end{equation}
The eigenvalues are $\lambda_1 = -1/x^2 < 0$ and $\lambda_2 = 0$.
Since $H_\psi$ is negative semidefinite for all $x > 0$, $\psi$ is concave. \qedhere
\end{proof}

\subsection{Spot Price}

\begin{proposition}[Spot Price from Trading Function]\label{prop:spot-price}
The spot price implied by the trading function is:
\begin{equation}\label{eq:spot-price}
  p = -\frac{\partial \psi / \partial x}{\partial \psi / \partial y}
    = -\frac{1/x - 1}{1}
    = 1 - \frac{1}{x}
    = \frac{1-x}{x}
\end{equation}
\end{proposition}

\begin{proof}
$\partial \psi / \partial x = 1/x - 1$ and $\partial \psi / \partial y = 1$.
The marginal rate of substitution is $-(1/x - 1)/1 = 1 - 1/x = (x-1)/x$.

Since $x \in (0,1]$, we have $x - 1 \leq 0$, so $(x-1)/x \leq 0$.
However, price must be non-negative. The sign convention depends on direction:
$p = |1/x - 1| = (1-x)/x$ for $x \in (0,1]$.

This is consistent with the reserve curve: $x = 1/(1+p) \implies p = (1-x)/x$. \qedhere
\end{proof}

\subsection{Canonical Form Analysis}

\begin{remark}[Non-Homogeneity]\label{rem:non-homogeneous}
$\psi(x,y)$ is \emph{not} 1-homogeneous:
\[
  \psi(\lambda x, \lambda y)
  = \lambda y + \ln(\lambda x) + 1 - \lambda x
  = \lambda y + \ln(\lambda) + \ln(x) + 1 - \lambda x
  \neq \lambda \psi(x, y)
\]
The $\ln(\lambda)$ term breaks homogeneity. This is expected for a fixed-supply CFMM
where the invariant curve is not scale-invariant. The Geometry paper
\cite{angeris2023geometry} establishes that canonical trading functions are
1-homogeneous; our $\psi$ is therefore not in canonical form but is a valid
(non-canonical) concave trading function.
\end{remark}

\subsection{Degeneracy Check}

\begin{remark}[Non-Degeneracy]\label{rem:non-degenerate}
The trading function $y = x - \ln(x) - 1$ is strictly convex in $x$
(since $y''(x) = 1/x^2 > 0$), meaning the reserve curve has non-zero curvature
everywhere. This is \emph{not} a degenerate (linear) CFMM.
The CFMM trades at a continuum of prices, not a single price point.
\end{remark}

\subsection{Liquidity Scaling}\label{subsec:liquidity-scaling}

With $L$ units of liquidity, reserves scale as $x \to L \cdot u$ where $u = x/L$
is the per-liquidity risky reserve. The trading function becomes:

\begin{equation}\label{eq:scaled-invariant}
  \frac{y}{L} = u - \ln(u) - 1
\end{equation}

\begin{proposition}[Pro-Rata Minting]\label{prop:pro-rata}
For an LP adding $\Delta L$ liquidity at price $p$, the required deposits are:
\begin{align}
  \Delta x &= \frac{\Delta L}{1 + p} \label{eq:mint-x}\\
  \Delta y &= \Delta L \cdot \Bigl(\ln(1+p) - \frac{p}{1+p}\Bigr) \label{eq:mint-y}
\end{align}
and the new shares satisfy $\Delta L / L = \Delta x / x = \Delta y / y$ (pro-rata).
\end{proposition}

\begin{proof}
At price $p$, the per-unit reserves are $x/L = 1/(1+p)$ and $y/L = \ln(1+p) - p/(1+p)$.
Scaling by $\Delta L$ gives the deposit amounts. The pro-rata property follows from
linearity of the reserve functions in $L$. \qedhere
\end{proof}

\subsection{Trading Set}

\begin{definition}[Trading Set]\label{def:trading-set}
The trading set is:
\begin{equation}\label{eq:trading-set}
  \TradingSet = \bigl\{(x, y) \in \mathbb{R}_{>0} \times \mathbb{R}_{\geq 0}
  : y + \ln(x) + 1 - x \geq 0 \bigr\}
\end{equation}
By \cref{prop:psi-concave}, $\psi$ is concave, so $\TradingSet$ is convex
(\cite[Theorem~1]{angeris2023geometry}).
\end{definition}
